\documentclass{article}
\usepackage{amsmath}
\usepackage{amssymb}
\usepackage{graphicx}
\usepackage{hyperref}
\usepackage{geometry}
\geometry{margin=1in}

\title{Closed form solution for the CatFlow optimization problem}
\author{Edgar Duc}
\date{\today}

\begin{document}
\maketitle

\section{One-dimensional case}
We denote as $\Delta^{K-1}$ the $(K-1)$-dimensional simplex, i.e. the set of $K$-dimensional vectors with non-negative entries that sum to 1. \\
We denote as $e_k$ the $k$-th standard basis vector in $\mathbb{R}^K$. \\
Let $\mathcal{D}_n = \{ x_i : 1 \leq i \leq n \} \subset \{e_1, \ldots, e_K\}$  denote the dataset.  \\
Let $\hat{p}_{data} = \frac{1}{n} \sum_{i=1}^n \delta_{x_i}$ denote the empirical distribution of the dataset. \\
Let $p_0 \in \mathcal{P}(\mathbb{R}^K)$ denote the distribution of the initial state $x_0$. \\

Let $X_1 \sim \hat{p}_{data}$, $X_0 \sim p_0$, $t \sim \mathcal{U}([0, 1])$.
Let $X_t = (1-t)X_0 + tX_1$.\\
We denote as $p_t$ the distribution of the state $x_t$ at time $t$. \\
Let $q : \mathbb{R}^K \times [0, 1] \to \Delta^{K-1}$ denote a velocity field. \\

The CatFlow loss is a cross-entropy loss, defined as:
\begin{equation}
    \mathcal{L}(q) = - \mathbb{E} \left[ \sum_{k=1}^{K} X_1^{(k)} \log q^{(k)}(X_t, t) \right]
\end{equation}
The optimization problem is to find the velocity field $q$ that minimizes the loss:
\begin{equation}
    \min_{q} \mathcal{L}(q)
\end{equation}

By conditioning on $X_t$ and $t$, we can rewrite the loss as:
\begin{equation}
    \mathcal{L}(q) = - \mathbb{E} \left[ \sum_{k=1}^{K} \mathbb{E}[X_1^{(k)} | X_t, t] \log q^{(k)}(X_t, t) \right]
\end{equation}

We know that $q_*$ minimizes the loss if and only if for all $x \in \mathbb{R}^K$ and $s \in [0, 1]$, we have:
\begin{equation}
    q_*(x, s) = \mathbb{E}[X_1 | X_t = x, t = s]
\end{equation}

We also know that :
\begin{equation}
    p_t(x | X_1 = x_i) = \frac{1}{(1-t)^{K}} p_0 \left( \frac{x - t x_i}{1-t} \right)
\end{equation}
Thus, by Bayes' rule, we have:
\begin{equation}
    \lambda_i(x, s) = \mathbb{P}[X_1 = x_i | X_t = x, t = s] = \frac{p_0 \left( \frac{x - s x_i}{1-s} \right) \hat{p}_{data}(x_i)}{\sum_{j=1}^n p_0 \left( \frac{x - s x_j}{1-s} \right) \hat{p}_{data}(x_j)} = \frac{p_0 \left( \frac{x - s x_i}{1-s} \right)}{\sum_{j=1}^n p_0 \left( \frac{x - s x_j}{1-s} \right)}
\end{equation}
Hence, we have:
\begin{equation}
\boxed{
    q_*(x,s)
    =
    \sum_{i=1}^{n} \lambda_i(x,s)\, x_i
}
\end{equation}

\section{Multi-dimensional case}
Let $D$ be the sequence length, and $K$ be the number of classes. \\
Let $\mathcal{D}_n = \{ x_i : 1 \leq i \leq n \} \subset \{e_1, \ldots, e_K\}^D$  denote the dataset.  \\
Let $\hat{p}_{data} = \frac{1}{n} \sum_{i=1}^n \delta_{x_i}$ denote the empirical distribution of the dataset. \\
Let $p_0 \in \mathcal{P}((\mathbb{R}^K)^D)$ denote the distribution of the initial state $x_0$. \\
Let $X_1 \sim \hat{p}_{data}$, $X_0 \sim p_0$, $t \sim \mathcal{U}([0, 1])$.
Let $X_t = (1-t)X_0 + tX_1$.\\
We denote as $p_t$ the distribution of the state $x_t$ at time $t$. \\
Let $q : ((\mathbb{R}^K)^D) \times [0, 1] \to (\Delta^{K-1})^D$ denote a velocity field. \\

The CatFlow loss is a cross-entropy loss \textbf{independently applied to each position in the sequence}, defined as:
\begin{equation}
    \mathcal{L}(q) = - \mathbb{E} \left[ \sum_{d=1}^{D} \sum_{k=1}^{K} X_1^{(d, k)} \log q^{(d, k)}(X_t, t) \right]
\end{equation}
The optimization problem is to find the velocity field $q$ that minimizes the loss:
\begin{equation}
    \min_{q} \mathcal{L}(q)
\end{equation}

As in the one-dimensional case, by conditioning on $X_t$ and $t$, we can rewrite the loss as:
\begin{equation}
    \mathcal{L}(q)
    =
    -\mathbb{E}\left[
        \sum_{d=1}^{D} \sum_{k=1}^{K}
        \mathbb{E}\left[X_1^{(d,k)} \mid X_t, t\right]
        \log q^{(d,k)}(X_t,t)
    \right]
\end{equation}

Since the loss is a sum of cross-entropies over positions, the minimizer is obtained pointwise by matching each categorical distribution to the corresponding conditional expectation. Therefore, for all $x \in ((\mathbb{R}^K)^D)$ and $s \in [0,1]$, we have:
\begin{equation}
    q_*(x,s)
    =
    \mathbb{E}[X_1 \mid X_t = x, t = s]
\end{equation}
or equivalently, for each position $d \in \{1,\dots,D\}$,
\begin{equation}
    q_*^{(d)}(x,s)
    =
    \mathbb{E}[X_1^{(d)} \mid X_t = x, t = s].
\end{equation}

Now let $x_i \in \{e_1,\dots,e_K\}^D$ denote the $i$-th sequence in the dataset. Since
\begin{equation}
    X_t = (1-t)X_0 + tX_1,
\end{equation}
conditioning on the event $X_1 = x_i$ gives
\begin{equation}
    X_t = (1-t)X_0 + t x_i.
\end{equation}
Hence, for $t \in [0,1)$,
\begin{equation}
    p_t(x \mid X_1 = x_i)
    =
    \frac{1}{(1-t)^{DK}}\,
    p_0\!\left(\frac{x - t x_i}{1-t}\right),
\end{equation}
where the factor $(1-t)^{-D(K-1)}$ comes from the affine change of variables on the ambient simplex product space, and is independent of $i$.

Therefore, by Bayes' rule,
\begin{equation}
    \lambda_i(x,s)
    :=
    \mathbb{P}[X_1 = x_i \mid X_t = x, t = s]
    =
    \frac{
        p_0\!\left(\frac{x - s x_i}{1-s}\right)\hat{p}_{data}(x_i)
    }{
        \sum_{j=1}^{n}
        p_0\!\left(\frac{x - s x_j}{1-s}\right)\hat{p}_{data}(x_j)
    }.
\end{equation}

Since $\hat{p}_{data}(x_i)=\frac{1}{n}$ for all $i$, this simplifies to
\begin{equation}
    \lambda_i(x,s)
    =
    \frac{
        p_0\!\left(\frac{x - s x_i}{1-s}\right)
    }{
        \sum_{j=1}^{n}
        p_0\!\left(\frac{x - s x_j}{1-s}\right)
    }.
\end{equation}

We finally obtain the closed-form solution:
\begin{equation}
\boxed{
    q_*(x,s)
    =
    \sum_{i=1}^{n} \lambda_i(x,s)\, x_i
}
\end{equation}

Equivalently, at each position $d$,
\begin{equation}
    q_*^{(d)}(x,s)
    =
    \sum_{i=1}^{n} \lambda_i(x,s)\, x_i^{(d)},
\end{equation}

Thus, exactly as in the one-dimensional case, the optimal predictor is a barycenter of the dataset points, with posterior weights given by the probability that each training sequence $x_i$ generated the observation $x$ at time $s$.\\
One could think that the solution would end up being a product of independent categorical distributions at each position, but this is only the case if the initial distribution $p_0$ and the data distribution $\hat{p}_{data}$ both factorize as a product of independent distributions at each position.

\end{document}